\section{Experimental Study}
\label{sec:experimental-setup}
We examine the path from a tested correction to an unassisted policy in a
development campaign and a subsequent one-update validation. The campaign
and validation are sequential diagnostics with implementation changes,
not randomized comparisons of repair representations.

\subsection{Setting and Comparisons}
\label{sec:experiment-design}
\paragraph{Tasks and policy.}
We use text-based ALFWorld \citep{shridhar2021alfworld} and the
Qwen2.5-3B-Instruct lineage \citep{qwen2025report}, updated with LoRA
\citep{hu2021lora}. The 3,553 training games are partitioned into 2,843
TRAIN\_UPDATE, 355 TRAIN\_SELECT, and 355 TRAIN\_AUDIT games. Only the
first pool supplies repair and learning evidence. Selection uses the
second pool. The final benchmark has 134 unseen and 140 seen games,
kept outside adaptive research and checkpoint selection.

\paragraph{Execution and statistical units.}
The policy receives the goal, current observation, eight executed
transitions, interface feedback, and the complete admissible menu. The
strict interface accepts one JSON action. Episodes allow 30 environment
steps and 60 policy attempts, with termination after three consecutive
nonexecuted attempts. Repair verification uses five paired seeds
(17, 31, 47, 73, 101); the final selection comparison uses one paired seed
(17) on all 355 tasks. Repetitions at one source are not independent
tasks. We report exact counts and descriptive differences for one trained
candidate, without claims of significance or training-seed robustness.
Appendix~\ref{app:configuration} records the executed configuration.

\paragraph{Scope of the evidence.}
We report completed intervention tests, native training labels, and paired
selection outcomes. Matched action-only, equal-compute, filter, hierarchy,
and memory ablations have not been completed for this update. The released
checkpoint benchmark supplies external context; it does not isolate
RePair's components. The RePair candidate has no completed final
seen/unseen benchmark result.

\subsection{Which Repairs Help the Current Policy?}
\label{sec:verification-results}
\begin{table}[!htbp]
\caption{\textbf{Completed same-state tests.} Counts are source-state
experiments; each has five pairs and ten branches. N denotes stable
neutral failure. Validations reuse R5 evidence and change the selected
portfolio and executable representation. They are not matched ablations.}
\label{tab:verification}\centering\small
\begin{tabular}{lrrrrr}\toprule
Run & States & Benefit & Harm & N & Uncertain\\\midrule
R1 & 11 & 0 & 0 & 11 & 0\\
R2 & 6 & 1 & 0 & 5 & 0\\
R3 & 2 & 0 & 0 & 2 & 0\\
R4 & 1 & 0 & 0 & 1 & 0\\
R5 & 1 & 0 & 0 & 1 & 0\\\midrule
Guarded-strategy validation & 17 & 0 & 0 & 16 & 1\\
Executable-strategy validation & 5 & 4 & 0 & 1 & 0\\\bottomrule
\end{tabular}
\end{table}

Table~\ref{tab:verification} separates the five development rounds from two
later validations that reuse R5 trajectories and Analyzer outputs. The
development rounds yield one stable Benefit among 21 tested states;
the other 20 are neutral failures. The first guarded-strategy validation
has no stable Benefit among 17 states, although two of its 85 pairs show
an F1-only success at the same source. The later executable-strategy
validation yields four stable Benefits among five selected states.
All five repetitions succeed in F1 and fail in F0 at each beneficial
state. The higher yield motivates examining what the successful programs do;
the sequentially selected portfolios leave the effect of each design change unresolved.

The trajectories explain why many valid corrections remain neutral.
In R3, reaching the fridge enables cooling an egg, but the continuation
then cools it 25 times instead of placing it. In a CD task, navigation
reveals the target, but the policy attempts placement three times without
first acquiring it. In R4, using a lamp is followed by 28 examinations of
the same side table. The intervention executes; the remaining dependency
chain fails. Appendix~\ref{app:behavior} records further cases.

The four beneficial programs in the later validation execute two, four,
six, and four actions, respectively, covering cleaning, placement,
heating-and-placement, and object-in-light tasks. All 20 successful F1
branches finish within the registered program, with zero subsequent
policy calls. The successful behavior therefore comes from the programs at these
selected states. Whether their verified actions become useful lessons
for the policy is the next question.

The heated-cup case makes the distinction concrete. From the restored
state, the program acquires the cup, reaches the microwave, heats it,
opens the microwave, reaches the cabinet, and places the cup. Completion
follows this six-action sequence without another policy decision. In
the remaining neutral case, the program locates and acquires a credit
card, but the three subsequent policy calls fail to execute and placement
remains unfinished. A useful intermediate state becomes a completed task
only when the remaining prerequisites are also resolved. These cases
locate the handoff that a trained policy would need to learn.

\subsection{What Enters Training?}
\label{sec:target-results}\label{sec:ablations}
The four Benefit states contribute four execution views and four strategy
views. The native-label audit counts 51 action-view and 2,546 strategy-view
loss-bearing tokens per pass. Four passes produce eight optimizer updates
and 10,388 supervised token exposures. Prompt and padding tokens are
masked. The current adapter teaches the initial action and associated strategy
at each source, leaving the remaining successful actions without
stepwise execution targets.
\begin{table}[!htbp]
\caption{\textbf{Verified program coverage and supervised action coverage.}
Each row is one Benefit source, with five successful F1 repetitions and
zero policy continuation calls. Action counts are per repetition; native
training counts are per source, not multiplied by the five seeds.}
\label{tab:coverage}\centering\small
\begin{tabularx}{\linewidth}{Xrrr}\toprule
Repair task & Program actions & Action targets & Strategy targets\\\midrule
Clean pan & 2 & 1 & 1\\
Place potato & 4 & 1 & 1\\
Heat and place cup & 6 & 1 & 1\\
Examine newspaper with lamp & 4 & 1 & 1\\\midrule
Total & 16 & 4 & 4\\\bottomrule
\end{tabularx}
\end{table}


For example, the heated-cup training pair supplies an action target for
acquiring the cup and a strategy description of the intended plan. The
five later action positions are exercised by the verifier but are not
separate action-prediction examples. The auxiliary view may communicate
the plan, while execution from the intervening observations remains a
generalization demand. Table~\ref{tab:coverage} makes this coverage gap
visible without treating each repetition as a new training source.

Thus verification coverage and training coverage differ: a program can
complete several prerequisites while the action target teaches only its
first decision. Strategy text supplies most supervised tokens; its gradient influence
also depends on token losses and accumulation-group composition
(Appendix~\ref{app:training-details}). The comparison below tests this
eight-example update as a whole. Isolating the strategy objective would
require an action-only update from the same four sources.

\subsection{What Remains After the Research Harness Is Removed?}
\label{sec:main-results}
\begin{table}[!htbp]
\caption{\textbf{Scaffold-free selection after the dual-view update.}
TRAIN\_SELECT only: 355 paired tasks, seed 17, one trained candidate.
Counts concern episodes except where marked as calls. Goal completion
is the mean terminal fraction of satisfied environment goal conditions.}
\label{tab:ablations}\centering
\begin{tabularx}{\linewidth}{Xrr}\toprule
Measure & Parent & Candidate\\\midrule
Task successes & 3 / 355 & 3 / 355\\
Mean goal-condition completion (\%) & 10.845 & 10.704\\
Consecutive-nonexecution terminations & 261 & 252\\
Environment-budget terminations & 91 & 100\\
Out-of-menu actions (calls) & 1,556 & 1,612\\
Format failures (calls) & 135 & 119\\\bottomrule
\end{tabularx}
\end{table}

On TRAIN\_SELECT, both parent and candidate solve 3 of 355 tasks (0.85\%).
Two successes are shared; one task improves and one regresses. Mean terminal
goal-condition completion decreases from 10.845\% to 10.704\%, so the
registered success-first, progress-on-ties rule retains the parent.
The comparison comprises 710 completed episodes with memory and
intervention disabled. It is a selection result, not a final benchmark.

Looking at task identities prevents the equal totals from hiding
redistribution: 351 tasks fail under both policies, two succeed under
both, and only two switch outcome. The secondary goal measure improves
on five tasks, worsens on six, and is unchanged on 344. We therefore find
neither a net completion gain nor a compensating aggregate progress gain
under this acceptance rule. The two discordant success outcomes support
an exact descriptive comparison, rather than a claim about training-seed
robustness.

Most failures concern execution rather than unavailable observations.
The parent and candidate terminate 261 and 252 tasks, respectively, after
consecutive nonexecuted attempts. Their traces contain 1,556 and 1,612
out-of-menu actions despite visible admissible menus. Fewer early stops
coexist with more budget-exhausted episodes (91 versus 100), leaving
terminal success unchanged. In a task requiring an alarm clock to be
examined with a desk lamp, both policies emit an unavailable examine
command even though the observation contains the clock and the menu
offers a take action. The missing step is therefore visible in the
interface; the policy fails to use it. This example links the aggregate
execution errors to a prerequisite failure of the same kind observed
after local repairs.

The existing reference-checkpoint benchmark is retained as task context
in Appendix~\ref{app:aggregation}; all policy-update conclusions here use
the paired TRAIN\_SELECT comparison.

\subsection{Research Cost and the Cross-Round Record}
\label{sec:discovery-results}
The five development rounds consume 210 completed verification branches
over 21 tested states. The guarded and executable-strategy validations
consume 170 and 50 branches, respectively. They reuse R5 collection and
analysis, so they are neither fresh formal rounds nor independent samples
of discovery yield. The last portfolio costs 50 branches for four Benefit
states, or 12.5 branches per Benefit, excluding upstream research and
earlier unsuccessful portfolios.

The campaign stops after five valid rounds under its three-round
no-promotion patience rule. One infrastructure-invalid attempt, repeated
engineering recoveries, and one user-authorized promotion exception remain
part of the record. The later validation completes training and selection
but encounters an ownership check during history closeout. We therefore
report its verified selection decision separately from campaign closure,
and make no claim of uninterrupted autonomous execution. Full API cost,
matched discovery efficiency, and further policy improvement remain
unmeasured in this study.
